\begin{frame}{Physically constrained correction}
\small

\begin{block}{Proposed model architecture: same model for all signals}

\large
\[
\boxed{y_\theta(s_1,s_2)=g_\theta(s_1)-g_\theta(s_2)}
\]

\small
\[
\theta^\star=\arg\min_\theta\frac{1}{N}\sum_{i=1}^{N}\left(y_\theta(s_{1}^i,s_{2}^i)-y^i_{\mathrm{true}}\right)^2,
\qquad
y^i_{\mathrm{true}}=\delta(s_{1}^i,s_{2}^i)+\varepsilon^i
=\Delta t_{\LED}^i-\widehat C_{1,2}-TOF
\]
\end{block}

\vspace{0.15em}

\begin{columns}[T,totalwidth=\textwidth]
\begin{column}{0.48\textwidth}
\begin{alertblock}{Corrected timing estimate}
\large
\[
\boxed{
\widehat{\Delta t}=\Delta t_{\LED}-\widehat C_{1,2}-y_\theta(s_1,s_2)
}
\]
\small
The final output is still a corrected arrival-time difference, so it can be evaluated with the same CTR metric.
\end{alertblock}
\end{column}

\begin{column}{0.48\textwidth}
\begin{block}{Why this is safer than generic pair-ML}
\begin{itemize}
  \item Same pulse scorer for every detector.
  \item Learns a timing-error correction, not an absolute position label.
  \item No need for multiple source positions: a single central source is sufficient.
\end{itemize}
\end{block}
\end{column}
\end{columns}

\vspace{0.25em}
\mainmsg{The proposed formulation gives a translation-invariant, antisymmetric,
and more general timing correction.}
\end{frame}
